% sec2b.tex — c05｜原文 p13–18（含 p19 顶端一句收尾）｜§2.2 规范理论的指标、§2.3 重组不变性、§2.4 拉氏量概览
% 编号公式 (2.17)–(2.29) 共 13 条；编号环境：习题 2.2、2.3；表 2（tabular 复刻）。
% 说明：§2.4 末句跨页到 p19 顶端，为保证句子完整，本块含 p19 开头一句；sec3a 从 §3 起。
% [原文 p.13]
\subsection{规范理论的指标}

除了 $\sN=1$ 手征多重态之外，四维超共形理论中随处可见的另一个多重态是 $\sN=1$ 矢量多重态。矢量多重态的指标可以用同样的方式计算，即先计算单字母指标，再取齐次指数。我们在表 2 中列出了 $\delta=0$ 的矢量多重态字母。
\begin{exercise}[2.2]
相对于 $\tilde Q_{\dot+}$ 计算矢量多重态的单字母指标。
\end{exercise}
\noindent\textbf{表 2.} 本表只列出 $\delta=0$ 的字母。这些是对指标有贡献的字母。

\begin{center}
\begin{tabular}{C{4.6em} C{1.4em} C{1.4em} C{1.4em} C{1.4em} C{1.4em} C{2.2em}}
\toprule
字母 & $\Delta$ & $j_1$ & $j_2$ & $r$ & $\delta$ & $\sI$ \\
\midrule
$\bar\lambda_{\dot+}$ & $\frac32$ & $0$ & $\frac12$ & $-1$ & $0$ & $-p$ \\
$\bar\lambda_{\dot-}$ & $\frac32$ & $0$ & $-\frac12$ & $-1$ & $0$ & $-q$ \\
$F_{++}$ & $2$ & $1$ & $0$ & $0$ & $0$ & $pq$ \\
$\partial^{\dot\alpha}_{+}\bar\lambda_{\dot\alpha}$ & $\frac52$ & $\frac12$ & $0$ & $-1$ & $0$ & $pq$ \\
$\partial_{+\dot+}$ & $1$ & $\frac12$ & $\frac12$ & $0$ & $0$ & $p$ \\
$\partial_{-\dot+}$ & $1$ & $-\frac12$ & $\frac12$ & $0$ & $0$ & $q$ \\
\bottomrule
\end{tabular}
\end{center}

因此，矢量多重态的单字母指标为
\[
i_V(a_i:p,q)=\frac{-p-q+2pq}{(1-p)(1-q)}\chi_{\adj}(a_i)
=\Big(-\frac{p}{1-p}-\frac{q}{1-q}\Big)\Big(\sum_\alpha a^\alpha+N\Big) .
\tag{2.17}
\]
这里我们已把单字母指标乘以规范群伴随表示的特征标。这是因为矢量多重态总处于伴随表示。伴随表示的特征标为 $\chi_{\adj}(a_i)=\sum_\alpha a^\alpha+N$，其中 $\alpha$ 是根向量。加性因子 $N$ 对应权向量为零的 Cartan 生成元。

既然已经知道手征多重态与矢量多重态的指标，我们就做好准备计算规范理论的指标了。依赖于物质内容，低能下的超对称规范理论流向非平庸的超共形场论。很可能无法精确计算它们的谱。然而超共形指标容易计算。诀窍在于意识到指标在规范耦合改变下不变。强耦合规范理论的指标可以通过把规范耦合设为零来计算。

% [原文 p.14]
这意味着我们只想枚举由 $\delta=0$ 的字母构造出的词。此类词上的配分函数就是单字母指标的齐次指数。这些就是四维中能构造出的全部超对称局域算符\footnote{在三维修与五维中，还有额外局域算符不能用拉氏量中基本场的字母构造。在三维它们称为「单极子算符」，在五维称为「瞬子算符」。计入它们意味着分别对 $S^2$ 与 $S^4$ 上非平凡磁通的分区求和。本讲义只关注四维理论。}。然而在规范理论中，我们只对构造规范变换下不变的算符感兴趣。这些算符是单态，即在规范群的平庸表示下变换。因此，规范理论的指标为
\[
\sI(b_k;p,q)=\frac{1}{|W|}\oint\Big(\prod_{i=1}^{N}\frac{\dd a_i}{2\pi\ii a_i}\Big)\Delta(a_i)\,\sI_V(a_i;p,q)\prod_{\phi_i}\sI_{\phi_i}(a_i,b_k;p,q) .
\tag{2.18}
\]
这里 $a_i$ 是规范对称性的逸度，$b_k$ 是手征多重态可能带有的所有其它对称性（整体对称性）的逸度。

利用 $\Delta(a_i)=(PE[\sum_\alpha a^\alpha])^{-1}$ 这一事实，我们可以化简组合
\[
\begin{aligned}
\Delta(a_i)\sI_V(a_i;p,q)&=PE\Big[\Big(-\frac{p}{1-p}-\frac{q}{1-q}\Big)\chi_{\adj}(a_i)-\sum_\alpha a^\alpha\Big]\\
&=PE\Big[\frac{pq-1}{(1-p)(1-q)}\Big(\sum_\alpha a^\alpha\Big)+N\Big(-\frac{p}{1-p}-\frac{q}{1-q}\Big)\Big]\\
&=\kappa^N\prod_\alpha\Gamma(pq\,a^\alpha),\qquad \kappa\equiv(p,p)(q,q),\quad (a,q)\equiv\prod_{i=0}^{\infty}(1-aq^i) .
\end{aligned}
\tag{2.19}
\]
函数 $(a,q)$ 称为 Pochhammer 符号。我们用简写 $\kappa$ 表示单个 $\U(1)$ 矢量多重态的指标。

利用这一化简，规范理论的指标可以完全用椭圆 gamma 函数（以及 Pochhammer 符号）表示出来。
\[
\begin{aligned}
\sI(b_k;p,q)&=\frac{\kappa^N}{|W|}\\
&\quad\times\oint\frac{\dd a_i}{(2\pi\ii)^N}\Gamma(pq\,a^\alpha)\prod_{\phi_i}\prod_{\rho\in R_i^G,\rho'\in R'^F_i}\Gamma\big((pq)^{\frac12 r_i}a^\rho b^{\rho'}\big) .
\end{aligned}
\tag{2.20}
\]
这里我们给手征多重态指定了任意的 R 荷 $r_i$。如前所述，它由超势确定，此外可能还要用 a-极大化。

\paragraph{例子} 现在我们给出一个 $\sN=1$ 规范理论中指标计算的例子，以说明超共形指标的威力。考虑 $\sN=1$、带三味基础与反基础夸克的 $\SU(2)$ 规范理论。该理论具有 $\SU(3)\times\SU(3)\times\U(1)$ 整体对称性。由于阿贝尔因子的存在，我们需要借助 a-极大化来确定正确的 R 荷。它们的结果为 $1/3$。
% [原文 p.15]
利用规范理论指标的一般公式 (2.20)，我们把该规范理论的指标写成积分
\[
\sI=\kappa\oint\frac{\dd z}{4\pi\ii z}\frac{1}{\Gamma(z^{\pm2})}\prod_{i=1}^{3}\Gamma\big((pq)^{\frac16}b\,u_i z^{\pm1}\big)\Gamma\big((pq)^{\frac16}b^{-1}v_i z^{\pm1}\big) .
\tag{2.21}
\]
这里 $\prod_{i=1}^3 u_i=\prod_{i=1}^3 v_i=1$，这些逸度参数化味对称性 $\SU(3)_u\times\SU(3)_v$——它转动基础与反基础夸克——而 $b$ 参数化重子的 $\U(1)_b$。通常，作为物理学家，我们会令 $p=xy$、$q=x/y$，把该积分按 $x$ 作 Taylor 级数展开来处理；但这个特定的积分已在数学文献中被研究过。引人注目的是，它允许闭式表达式
\[
\sI=\prod_{i<j}\Gamma\big((pq)^{\frac13}t_it_j\big) .
\tag{2.22}
\]
这里 $\{t_i\}=\{bu_i,b^{-1}v_i\}$。这称为 Spiridonov 的 beta 函数恒等式 [19]。
\begin{exercise}[2.3]
在 $x$ 的 Taylor 级数展开中验证恒等式 (2.22)。
\end{exercise}
这一形式极具启发性。它是十五个手征多重态的指标，这些多重态与一个超势相互作用，该超势赋予每个多重态 R 荷 $2/3$。此外，$u$ 与 $v$ 逸度组合成六个 $t_i$，使得指标在它们的置换下不变，这暗示该理论具有 $\SU(6)$ 对称性，其逸度为 $t_i$（对称性是 $\SU(6)$ 而非 $\U(6)$，因为 $\prod_i t_i=1$）。这样的超势正是由这十五个手征场构造出的 $6\times6$ 反对称矩阵的 Pfaffian。

这暗示规范理论与仅含手征多重态的理论之间存在对偶。这确实成立。这是众多著名的 Seiberg 对偶之一 [20]。碰巧这一特定对偶对超共形指标的推论已在数学文献中被独立研究过。但更一般地，以这种方式，所有 Seiberg 对偶都成为非平凡数学积分恒等式的来源\footnote{这些理论在大 $N$ 极限下的超共形指标已在 [21] 中研究。}。事实上，这类积分（称为椭圆超几何积分）已被数学家单独研究。参见 [22] 的导引与文献。更重要的是，从物理学家的观点看，超共形指标的匹配为非凡的非微扰对偶猜想提供了非平凡的检验。

\subsection{重组不变性}

在超对称量子力学中，Hilbert 空间构成超对称代数 $\SU(1|1)$ 的一个表示。有两类不可约表示。一个二维长多重态 $\{|\chi\rangle,|\psi\rangle\}$
\[
Q|\chi\rangle=|\psi\rangle,\qquad Q|\psi\rangle=0
\tag{2.23}
\]
% [原文 p.16]
以及一个一维短多重态 $\{|\chi\rangle\}$
\[
Q|\phi\rangle=0 .
\tag{2.24}
\]
仅从代数 $\{Q,Q^\dagger\}=H$ 就容易看出 $H|\phi\rangle=0$。若对理论作光滑的超对称形变，长多重态的能量可能降到零，在这一点它分裂成两个短多重态。反之，若在零能处有两个短多重态，则光滑的超对称形变一般会使它们重组形成一个长多重态。Witten 指标的关键性质是它在重组下不变，这正是它在光滑超对称形变下受保护的原因。

人们可以对超共形指标提出同样的表示论问题，并把它理解为一个在超共形代数所允许的一切重组下不变的对象。Lorenz Eberhardt 已经讲过超共形代数及其长、短表示。我们将为此借鉴那一讨论。虽然重组规则最早在 [23] 中讨论，我们将遵循 [12] 的叙述。让我们回忆，

三类缩短条件
\[
\begin{aligned}
A_1:\ \epsilon^{\alpha\beta}\mathcal{Q}_\alpha([j_1,j_2]^r_\Delta)_\beta=0
&\quad\Rightarrow\quad \Delta=2+2j_1-\frac32 r
&\qquad\text{零态：}\quad &[j_1-\tfrac12,j_2]^{r-1}_{\Delta+\frac12}\\
A_2:\ \epsilon^{\alpha\beta}\mathcal{Q}_\alpha\mathcal{Q}_\beta[0,j_2]^r_\Delta=0
&\quad\Rightarrow\quad \Delta=2-\frac32 r
&\qquad\text{零态：}\quad &[0,j_2]^{r-2}_{\Delta+1}\\
B_1:\ \mathcal{Q}_\alpha[0,j_2]^r_\Delta=0
&\quad\Rightarrow\quad \Delta=-\frac32 r
&\qquad\text{零态：}\quad &[\tfrac12,j_2]^{r-1}_{\Delta+\frac12}
\end{aligned}
\tag{2.25}
\]
我们采用了 [24] 中分类缩短条件的约定。考虑 $A_1$ 型缩短条件：当超共形多重态的原初态饱和幺正性界 $\Delta=2+2j_1-\frac32 r$ 时，态 $[j_1-\frac12,j_2]^{r-1}_{\Delta+\frac12}$ 变为零态。有趣的是，容易验证该零态也处在幺正性界上。如果你现在把带零态的那个多重态视为以零态为原初态的多重态，那么这个多重态也有零态，并且同样处在幺正性界上。在这一过程中 $j_1$ 的值不断减小，链条在 $j_1=0$ 处终止。在那种情形下，零态满足 $B_1$ 型缩短条件。总之，我们得到一列短多重态，使得任意两个相邻多重态都能重组形成一个长多重态，
\[
[j_1,j_2]^r_\Delta\to[j_1-\tfrac12,j_2]^{r-1}_{\Delta+\frac12}\to\ldots\to[0,j_2]^{r-2j_1}_{\Delta+j_1}\to[0,j_2]^{r-2j-2}_{\Delta+j_1+1} .
\tag{2.26}
\]
链条两端会发生什么？显然，链条不能任意向右延伸，因为共形维数在那个方向上递减。另一方面，右端的 $B_1$ 型多重态具有零态 $[\frac12,j_2]^{r-2j-3}_{\Delta+j+\frac12}$。这个零态有 $j_1\neq0$，并且容易验证它位于幺正性界之下，因此根本谈不上以它作为原初态构造多重态。这样，重组链在两个方向上都终止。

% [原文 p.17]
为着重组的用途，把末端的 $B_1$ 型多重态看作 $j_1=-\frac12$ 的 $A_1$ 型多重态是方便的，即
\[
[0,j_2]^{r}_{\Delta=-\frac32 r}\equiv[-\tfrac12,j_2]^{r+1}_{\Delta=-\frac32 r-\frac12} .
\tag{2.27}
\]
超共形指标需要捕捉模掉重组之后的全部谱信息；由于上述链条中任意两个相邻多重态都能重组，指标应当计入
\[
N_{r'=r-2j_1,j_2}=\sum_{p=-1}^{k}(-1)^p\,n\big([\tfrac{p}{2},j_2]^{r'+p}_{\Delta=2+p-\frac32(r'+p)}\big) .
\tag{2.28}
\]
确实，$A_1$ 型多重态 $[j_1,j_2]^r_\Delta$ 的指标正是 [25]，
\[
\sI_{[j_1,j_2]^r_\Delta}=(-)^{2(j_1+j_2)+1}\frac{(pq)^{\frac12(-r+2j_1+2)}\chi_{j_2}\big(\sqrt{q/p}\big)}{(1-p)(1-q)} .
\tag{2.29}
\]
若我们在重组链中的一族多重态上求值，它恰好计入 $N_{r'=r-2j_1,j_2}$。这是因为，除符号交替之外，链条中所有多重态的指标都相同。

$\sN=1$ 超共形代数为 $\SU(2,2|1)$，该代数中单个超荷的对易子为 $\SU(2,1)$。指标必须是该代数的特征标。确实，该代数的两个 Cartan 生成元在 $\SU(2,2|1)$ 中的嵌入为 $r-2j_1$ 与 $j_2$。这与重组规则的预期一致。

\subsection{中场休息：超对称拉氏量概览}

\paragraph{$\sN=0$ 拉氏量：} 一条经验法则是：超对称越多，指定拉氏量所需的数据越少。例如，若想构造标量、费米子与规范场的普遍拉氏量。首先必须指定物质场在规范对称性下的荷（或更一般地，表示）。除规范相互作用外，物质还可以带有自身的相关相互作用，例如 Yukawa 项与四次标量项，并以保持规范对称性的方式出现。最后，费米子的数目及其表示需要仔细选取，以使得规范对称性无反常。

\paragraph{$\sN=1$ 拉氏量：} 有了 $\sN=1$ 超对称，分类多少简化了一些（尽管远不及 $\sN=2$ 那样简单）。标量与费米子在一个多重态中，即手征多重态；规范场与规范微子在一个多重态中，即矢量多重态。指定手征多重态在规范对称性下荷的表示，就完全固定了手征—矢量相互作用。对 $\sN=0$，物质相互作用有两类且彼此独立；但有了 $\sN=1$ 超对称，它们彼此关联，并可写成一个超势 $W(\Phi_i)$。同样，需要确保物质内容使得规范对称性无反常。

% [原文 p.18]
\paragraph{$\sN=2$ 拉氏量：} $\sN=2$ 拉氏量用 $\sN=2$ 超多重态（由一对手征多重态组成）与 $\sN=2$ 矢量多重态（由 $\sN=1$ 矢量多重态与 $\sN=1$ 手征多重态组成，二者都在规范群的伴随表示下变换）来构造。具有 $\sN=2$ 超对称的拉氏量比 $\sN=1$ 的容易指定得多。只需指定矢量—物质的耦合；物质—物质之间的耦合被 $\sN=2$ 超对称禁止。前者只需通过指定超多重态在规范群下的表示来给出。这使得可以用称为箭图的便利图解记号来描述 $\sN=2$ 理论。此外，由于物质成对地处于共轭表示中，$\sN=2$ 理论的规范对称性总是无反常的。因此以这种方式指定的拉氏量总是好的拉氏量（不过它可能不是渐近自由的）。当规范耦合的单圈 beta 函数为零时，规范耦合充当恰当边缘形变。在那种情形下，相互作用超共形理论由复化规范耦合参数化。

\paragraph{$\sN=3$ 拉氏量：} 不存在不具有 $\sN=4$ 对称性的 $\sN=3$ 拉氏量。不过近年来发现了一些非拉氏量的 $\sN=3$（但非 $\sN=4$）理论 [26–28]。$\sN=1$ 在 $\sN=3$ 中的对易子秩为二。因此可以再开启两个逸度，使指标成为四个变量的函数。关于 $\sN=3$ 理论指标的讨论见 [29]。

\paragraph{$\sN=4$ 拉氏量：} $\sN=4$ 拉氏量由伴随表示中的单个矢量多重态构造。它总是共形的，复化规范耦合参数化超共形理论的空间。$\sN=4$ 矢量多重态由伴随表示中的 $\sN=2$ 矢量多重态与 $\sN=2$ 超多重态组成。用 $\sN=1$ 多重态来说，它由一个 $\sN=1$ 矢量多重态与三个在伴随表示中变换的 $\sN=1$ 手征多重态组成。$\sN=1$ 在 $\sN=4$ 中的对易子秩为二\footnote{它并非人们天真猜测的三，原因在于 R 对称性是 $\SU(4)$ 而非 $\SU(4)\times\U(1)$。对 $\sN=1,2,3$，R 对称性是 $\SU(N)\times\U(1)$。}。因此可以再开启两个逸度，使指标成为四个变量的函数。

如果超共形代数的秩为 $r$，则最一般的指标是 $r-2$ 个逸度的函数。这是因为指标只计入 $\delta=0$ 的态，而我们只能对与 $Q$ 对易的荷开启逸度。
% [原文 p.19]
这两个条件把变量数减少 2。

% [Tutor Note: 术语：表 2 中的「字母」是单字母，$\sI$ 列给出该字母对单字母指标分子（或分母塔）的贡献，而不是完整指标。]
% [Tutor Note: 易错：表 2 只列 $\delta=0$ 的字母，这与表 1 不同；$\delta\neq0$ 的字母对指标无贡献，故不列出。]
% [Tutor Note: 易错：$\chi_{\adj}$ 中的加性 $N$ 来自 Cartan 生成元（零权），不要漏掉，否则 $i_V$ 与 (2.17) 不符。]
% [Tutor Note: 动机：指标在规范耦合变化下不变，因此可在弱耦合（$g=0$）下计算，再解析延拓到强耦合——这是指标能给出非平凡检验的根本原因。]
% [Tutor Note: 易错：(2.18) 是被积函数在规范群上的投影；$1/|W|$ 与 $\Delta(a_i)$ 一起实现 Weyl 投影，$\Delta$ 的幂次是 $+1$。]
% [Tutor Note: 易错：(2.20) 原文照录，但测度应为 $\prod_i \dd a_i/(2\pi\ii)$，且单独的 $\Gamma(pqa^\alpha)$ 应为 $\prod_\alpha\Gamma(pqa^\alpha)$；已登记 ERRATA。]
% [Tutor Note: 术语：Spiridonov beta 函数恒等式 (2.22) 即数学中的椭圆 beta 积分，$\SU(2)$ 三味指标的闭式正是它。]
% [Tutor Note: 动机：(2.21) 中 $z^{\pm1}$ 是简写 $\{z,z^{-1}\}$，$z^{\pm2}=\{z^2,z^{-2}\}$，不要漏掉任何一支。]
% [Tutor Note: 易错：Seiberg 对偶两边指标相等，这是对偶的非平凡检验；但它不证明对偶，只检验 BPS 谱。]
% [Tutor Note: 术语：长/短多重态、重组、零态、幺正性界——短多重态在重组下可与长多重态互换，指标对重组不敏感。]
% [Tutor Note: 易错：(2.24) 原文正文说短多重态是 $\{|\chi\rangle\}$ 而方程写 $Q|\phi\rangle=0$，符号不一致；已登记 ERRATA。]
% [Tutor Note: 易错：(2.26) 末项指数原文作 $r-2j-2$、下段作 $r-2j-3$，其中 $j$ 应为 $j_1$；译文照录，已登记 ERRATA。]
% [Tutor Note: 动机：(2.28) 的 $(-1)^p$ 正是长短多重态交替的符号，求和后把链条的贡献约化为一个数。]
% [Tutor Note: 术语：$\sN=1$ 指标是 $\SU(2,1)$ 的特征标，两个 Cartan 为 $r-2j_1$、$j_2$；这解释了 (2.29) 的形式。]
% [Tutor Note: 文献：[25] 给出 $A_1$ 型多重态指标的显式形式；[23] 最早讨论重组规则；[12] 为本文采用的叙述。]
% [Tutor Note: 动机：$\sN=2,3,4$ 的指标分别依赖 3、4、4 个变量，因为对易子的秩随超对称增加而增加。]
